\documentclass[conference]{IEEEtran}
\IEEEoverridecommandlockouts
\usepackage{cite}
\usepackage{amsmath,amssymb,amsfonts}
\usepackage{booktabs}
\usepackage{array}
\usepackage{url}
\usepackage{graphicx}
\usepackage{textcomp}
\usepackage{xcolor}
\def\BibTeX{{\rm B\kern-.05em{\sc i\kern-.025em b}\kern-.08em
    T\kern-.1667em\lower.7ex\hbox{E}\kern-.125emX}}

\begin{document}

\title{Why Do We Need Sleep? A Sleep Function Causal Atlas for Memory, Maintenance, and Immune Resilience}

\author{\IEEEauthorblockN{Four-Agent Research Pipeline}
\IEEEauthorblockA{Survey--Idea--ExperimentDesign--Author Workflow\\
Research proposal; no experiments executed}}

\maketitle

\begin{abstract}
Sleep occupies roughly a quarter to a third of human life, yet its scientific function is not captured by a single explanation. This paper proposes the Sleep Function Causal Atlas (SFCA), a design-first framework for testing how homeostatic sleep pressure, circadian phase, continuity, stage microstructure, and individual physiological state jointly produce domain-specific benefits. The proposal integrates memory and plasticity, brain maintenance and clearance-related physiology, immune defense, and systemic health without reducing them to one biomarker or one nightly duration. Its central hypothesis is that a decomposed, stage- and timing-aware sleep state predicts distinct outcomes better than duration alone, and that safe, approved perturbations can reveal shared functions and dissociable mechanisms. We formalize a two-process sleep state model, a hierarchical outcome model, domain-specific memory and immune endpoints, causal contrasts, mediation criteria, and a preference-explicit cross-domain utility layer. A randomized, counterbalanced, repeated-measures protocol separates prior wake from circadian timing and tests recovery, with controls for alertness, encoding opportunity, practice, stress, illness, medication, and measurement quality. The study remains design-only: no sleep restriction, circadian displacement, physiological measurement, or clinical intervention was executed. SFCA offers a direct answer to the motivating question: sleep is needed not for one universal purpose but because regulated sleep states coordinate multiple forms of neural plasticity, cellular maintenance, immune regulation, and bodily recovery whose relative value depends on time and individual state.
\end{abstract}

\begin{IEEEkeywords}
sleep, circadian rhythm, homeostatic sleep pressure, memory consolidation, neuroimmune regulation, causal inference, personalized physiology
\end{IEEEkeywords}

\section{Introduction}
Humans spend approximately 25--30 percent of their lives asleep. The time cost can make sleep appear passive, but sleep is an actively regulated physiological state. Sleep pressure accumulates during waking, the circadian clock gates the timing of sleep and wake, and sleep architecture changes across development and health \cite{BorbeLy1982,NHLBI2022}. Insufficient, mistimed, fragmented, or poor-quality sleep can impair learning, memory, attention, reaction, mood, immune defense, and long-term physical health \cite{NHLBI2022}. The scientific problem is therefore not whether sleep matters. It is how its distinct components produce distinct forms of benefit, and why the cost of disrupting one component differs across people and outcomes.

The phrase ``why do we need sleep?'' invites a single-purpose answer: memory consolidation, energy conservation, brain clearance, synaptic reset, immune repair, or evolutionary protection. Each hypothesis captures part of the evidence, but a single-purpose theory predicts too little. A person can show preserved item recognition but impaired temporal order; sleep loss can alter inflammatory signaling without an immediate memory deficit; a circadian shift can impair performance even when time in bed is matched; and a recovery night can reverse some effects while leaving others. The same intervention can thus reveal both shared and dissociable functions.

This paper introduces the Sleep Function Causal Atlas (SFCA), a research program selected through a four-agent workflow. SFCA treats sleep as a vector of interacting features rather than a scalar duration. The atlas maps these features to declared phenotypes in memory, neural plasticity, brain maintenance, immune defense, and systemic regulation. Its objective is not to accumulate correlations under the label of sleep quality. It is to identify which feature is changed, which endpoint responds, which alternative explanation is excluded, and for whom the response occurs.

The proposal makes four contributions. First, it defines a sleep phenotype that separates duration, timing, continuity, stage composition, prior wake, and physiological state. Second, it specifies a factorial repeated-measures design that makes homeostatic pressure and circadian phase separately testable. Third, it defines content-sensitive memory endpoints and multi-marker immune and maintenance endpoints within a hierarchical causal model. Fourth, it adds individual-response and safety gates so that personalization cannot silently trade memory against immune function, mood, or safety. The artifact is a research plan, not an execution report: no human participant was recruited, no sleep condition was imposed, and no physiological result is claimed.

\section{Background, Survey, and Research Gap}
\subsection{Sleep is a regulated state}
The two-process model provides a useful starting point. Process S, or homeostatic sleep pressure, rises during wake and dissipates during sleep. Process C, the circadian component, oscillates near 24 hours and gates the propensity for sleep and wake. The two processes interact: a person may be highly sleepy at an unusual clock time after prolonged wake, or alert at a usual bedtime when sleep pressure is low. Contemporary physiology adds orexinergic, adenosine, melatonin, cortisol, thermoregulatory, and environmental pathways, but the conceptual separation remains valuable because duration, prior wake, and circadian phase are often confounded in ordinary life.

The National Heart, Lung, and Blood Institute defines sleep deficiency broadly to include insufficient sleep, sleeping at the wrong time, poor-quality or incomplete sleep, and sleep disorders \cite{NHLBI2022}. Its public guidance identifies an internal body clock, a recurring circadian rhythm, and a sleep drive that builds across waking. It also reports that ongoing sleep deficiency can affect learning, memory, decision making, immune defense, cardiovascular health, metabolism, mood, and safety. These statements establish sleep as a multi-domain health requirement, but they do not establish the relative causal contribution of each sleep feature.

\subsection{Memory, plasticity, and replay}
Sleep and memory interact at multiple stages. Newly encoded memories may be stabilized, transformed, integrated with older knowledge, or selectively weakened during subsequent sleep. Slow oscillations, sleep spindles, hippocampal replay, neuromodulatory state, and rapid-eye-movement sleep have all been proposed as components of this process \cite{Walker2006,RaschBorn2013,Klinzing2019}. These mechanisms make different predictions. Replay predicts content-specific reactivation and improved relational integration. A synaptic-homeostasis account predicts broad changes in network gain and signal-to-noise. A circadian-gating account predicts performance changes at matched sleep duration but different phases. A recovery account predicts interactions with prior wake and repeated learning.

Memory must be measured as more than a single recall percentage. An item can be recognized while its source is confused; a sequence can be remembered without precise item identity; and a relation can be inferred from gist rather than replayed from the original event. SFCA therefore uses item identity, temporal order, source context, relational structure, confidence, response time, interference resistance, and alertness. The central memory outcome is a declared phenotype, not a claim to measure every autobiographical experience.

\subsection{Brain maintenance and clearance}
Sleep is also associated with cellular and fluid-transport processes that may support brain maintenance. Experimental work has linked sleep to increased clearance of metabolites in the mouse brain, motivating the glymphatic and interstitial-fluid literature \cite{Xie2013,Nedergaard2020}. This is an important mechanistic opportunity, but it is not a complete explanation of sleep. Measurement differs across species and modalities, and a change in a clearance-related marker does not by itself establish a cognitive or clinical benefit. SFCA treats clearance as a candidate mediator whose reliability, temporal ordering, and incremental explanatory value must be tested.

\subsection{Immune and systemic regulation}
Sleep and immunity are coupled through bidirectional signaling. Sleep loss and fragmentation can alter inflammatory pathways, leukocyte trafficking, endocrine regulation, and responses to infectious challenge, while inflammation can in turn alter sleep architecture \cite{Besedovsky2019,Irwin2015}. Sleep also supports cardiovascular, metabolic, endocrine, growth, and tissue-repair functions. Because illness, stress, medication, activity, diet, and socioeconomic timing affect both sleep and health, observational associations require careful causal interpretation. A robust atlas should measure a panel of immune and systemic outcomes, report heterogeneity, and preserve adverse or null outcomes.

\subsection{Survey gap ledger}
The Survey Agent accepted seven research gaps. G1 is the absence of a unified stage- and timing-aware operational definition of sleep benefit. G2 concerns the causal separation of homeostatic pressure and circadian phase. G3 concerns the distinction among replay, plasticity, synaptic renormalization, arousal, and encoding opportunity. G4 concerns the causal relationship between clearance-related physiology and cognitive or disease outcomes. G5 concerns individualized immune and systemic responses. G6 establishes that duration is insufficient without continuity, stage composition, timing, and prior wake. G7 concerns biomarkers that predict sleep resilience or response to intervention. Cross-domain utility and causal sleep-resilience phenotypes were retained as candidate rather than established gaps.

The evidence policy is bounded. NIH material supports claims about regulated sleep and health associations; sleep and memory reviews support mechanistic alternatives; clearance studies support a candidate maintenance pathway; and neuroimmune reviews support bidirectional sleep--immune coupling. None supports the claim that one mechanism is the sole purpose of sleep. The evidence registry is summarized in Table~\ref{tab:evidence} and frozen in the Survey manifest.

\begin{table}[t]
\caption{Evidence registry and permitted interpretation}
\label{tab:evidence}
\centering
\footnotesize
\setlength{\tabcolsep}{2pt}
\begin{tabular}{p{0.13\columnwidth}p{0.35\columnwidth}p{0.40\columnwidth}}
\toprule
Sources & Evidence family & Permitted interpretation \\
\midrule
E1,E2 & Regulation and health & Sleep is regulated by interacting pressure and clock processes; deficiency affects multiple domains. \\
E3,E4 & Memory and plasticity & Sleep can support memory processing; mechanisms and memory domains remain separable questions. \\
E5,E6 & Brain maintenance & Sleep-associated clearance is a candidate physiological pathway, not a universal explanation. \\
E7,E8 & Immune regulation & Sleep and immune signaling interact; confounding and individual variation require causal controls. \\
E9 & Individual variation & Age, chronotype, health, and state can moderate sleep timing and response. \\
\bottomrule
\end{tabular}
\end{table}

\section{Research Questions and Planned Contributions}
The primary question is: which interacting features of sleep are causally necessary for specific cognitive and physiological outcomes, and how do their effects depend on stage, timing, prior wake, and individual state?

This question is decomposed into five tests:
\begin{enumerate}
\item Does a decomposed sleep-state model predict memory, immune, maintenance, and systemic outcomes better than total duration?
\item Can factorial timing and prior-wake manipulations separate homeostatic pressure from circadian phase?
\item Which sleep features predict item, order, source, and relational memory beyond alertness and encoding opportunity?
\item Do clearance-related and immune markers mediate identifiable sleep effects, or do they provide domain-specific dissociations?
\item Can an individual-response model select a safe, endpoint-specific recovery schedule without assuming a single universal sleep need?
\end{enumerate}

The planned contribution is the SFCA atlas: a versioned map of sleep features, intervention contrasts, outcome domains, individual moderators, and failure modes. It combines a mechanistic layer, an outcome layer, a causal layer, and a decision layer. The mechanistic layer estimates pressure, phase, stages, and continuity. The outcome layer stores domain-specific phenotypes. The causal layer records interventions and controls. The decision layer promotes only findings that generalize beyond confounds and remain within safety bounds. This structure can support targeted sleep medicine, adaptive scheduling, and mechanistic neuroscience while keeping the claim proportional to the endpoint tested.

\section{Idea Source Checkpoints and Direction Selection Audit}
The Idea Agent explored four routes. R1 was a two-process perturbation matrix. R2 was a stage-resolved memory replay atlas. R3 was a neuroimmune maintenance bridge. R4 was a personalized sleep-resilience controller. R5, SFCA, combined the strongest components under explicit domain and safety gates.

R1 offers strong causal leverage but risks narrowing sleep to timing. R2 offers a sharp neural mechanism but cannot explain immune or systemic outcomes alone. R3 connects clinically meaningful domains but risks becoming a large correlational model. R4 is translationally attractive but can optimize a proxy without understanding why it works. SFCA was selected because it retains all four routes while forcing each to answer a distinct question. A memory improvement is not automatically immune recovery; a clearance marker is not automatically disease prevention; and a personalized schedule is not successful if it improves one domain while harming another.

The search rejected three attractive but scientifically weak directions. A memory-only theory cannot account for systemic and immune evidence. A clearance-only theory treats a candidate mechanism as a universal purpose. A duration-only predictor discards circadian timing, sleep stages, continuity, prior sleep debt, and individual physiology. The selected direction is ambitious because it seeks a unified causal atlas, but it is testable because it does not require all domains to move together.

The MCTS-style evolution checkpoints were: (N0) replace scalar sleep duration with a decomposed sleep phenotype; (N1) add memory, immune, maintenance, and systemic endpoints; (N2) add factorial timing and recovery to separate the two processes; and (N3) add hierarchical individual response and explicit safety utility. These edits turn a broad biological question into a portfolio of falsifiable hypotheses.

\section{Problem Definition, Assumptions, and Hypotheses}
\subsection{Definitions}
SFCA defines sleep need relative to a declared phenotype and time horizon. A short-term cognitive need is the sleep state required to preserve a specified memory or attention endpoint. A physiological need is the state required to maintain a declared immune, endocrine, metabolic, cardiovascular, or maintenance endpoint. A long-term health need concerns repeated exposure and risk, not one morning's performance. These definitions prevent the phrase “need sleep” from hiding different causal questions.

The sleep state vector separates duration $D$, continuity and quality $Q$, homeostatic pressure $H$, circadian phase $C$, stage features such as slow-wave activity and spindle density, rapid-eye-movement timing, and individual physiological state $U$. The outcome vector separates memory phenotypes, alertness, immune markers, clearance-related measures, systemic physiology, mood, and safety. A shared latent state may be proposed, but domain-specific outcomes are primary.

\subsection{Assumptions}
The design assumes that sleep features and outcomes can be measured with sufficient reliability under an approved protocol. It assumes that prior wake, light exposure, caffeine, medication, illness, activity, stress, and task practice can be recorded or constrained. It assumes that randomized crossover and recovery periods can reduce between-person confounding and carryover. It does not assume that one laboratory night represents habitual sleep, that a wearable fully identifies sleep stages, that a clearance marker is clinically decisive, or that a neural correlate is a complete mechanism.

\subsection{Hypotheses}
H1 states that decomposed sleep-state models outperform duration-only models on held-out cognitive, immune, maintenance, and systemic outcomes. H2 states that prior wake and circadian phase have separable causal effects when time in bed is matched. H3 states that stage microstructure predicts different memory components: replay-related measures will be more closely associated with content and relational retention, while broad alertness and pressure will affect response speed and encoding opportunity. H4 states that sleep-related immune and maintenance effects are partially dissociable from memory effects and vary across individuals. H5 states that a personalized, endpoint-specific controller improves recovery relative to a safe duration-matched schedule without increasing adverse outcomes. H6 predicts that confidence or subjective refreshment alone will not reliably identify physiological recovery.

\section{Study Design and Methods}
\subsection{D0: Scope, safety, and preregistration}
Before recruitment, preregister the population, age range, sleep-disorder screening, medical and psychiatric exclusion rules, maximum sleep loss, recovery schedule, stop rules, compensation, consent language, task order, randomization, endpoint hierarchy, power assumptions, missing-data plan, and adverse-event process. Participants at risk from altered sleep schedules must be excluded or referred for expert review. No participant may drive, operate hazardous equipment, or make high-stakes decisions during a sleep-loss condition.

The proposed work is a controlled laboratory or closely monitored digital protocol after institutional approval. It is not a recommendation for unsupervised sleep deprivation. Sleep medicine, clinical safety, immunology, statistics, privacy, and ethics reviewers must approve the protocol before any physical measurement. The present document contains no observed intervention or outcome.

\subsection{D1: Factorial two-process protocol}
Use a randomized, counterbalanced, within-person crossover design with washout and recovery nights. Core conditions are adequate sleep opportunity at habitual phase, modest safety-bounded sleep restriction at habitual phase, circadian displacement with matched time in bed, and recovery sleep. If safe and approved, a continuity-fragmentation condition separates duration from fragmentation. Light exposure, caffeine, meals, activity, medication, and task schedules are logged or standardized.

The operational homeostatic model is
\begin{equation}
H_{t+1}=\rho H_t+\alpha W_t-\beta S_t+\eta_t,
\label{eq:homeostatic}
\end{equation}
where $W_t$ and $S_t$ quantify wake and sleep exposure, $\rho$ captures persistence, and $\eta_t$ is unobserved disturbance. Circadian phase is represented by
\begin{equation}
C_{t+1}=C_t+\omega+\gamma L_t+\xi_t \pmod {2\pi},
\label{eq:circadian}
\end{equation}
where $L_t$ is measured light input, $\omega$ is intrinsic phase advance, and $\xi_t$ is phase uncertainty. Equations~\eqref{eq:homeostatic} and \eqref{eq:circadian} are operational models for separating experimental variables, not complete descriptions of sleep biology.

\subsection{D2: Sleep phenotype and measurements}
Measure sleep duration, onset and offset, efficiency, wake after sleep onset, fragmentation, stage composition, slow-wave activity, spindle density, rapid-eye-movement timing, circadian phase proxy, subjective sleepiness, alertness, and prior sleep debt. Use polysomnography or validated multimodal sensing when appropriate. Record light, activity, caffeine, alcohol, medication, illness symptoms, chronotype, stress, and engagement.

Let the measured sleep state be
\begin{equation}
z_t=[H_t,C_t,D_t,Q_t,\mathrm{SWA}_t,\mathrm{Spindle}_t,\mathrm{REM}_t,U_t]^\top,
\label{eq:sleepstate}
\end{equation}
where $Q_t$ denotes continuity and quality and $U_t$ denotes individual physiological state. Compare duration-only, timing-only, stage-only, and full decomposed-state models using held-out participant-session data. The comparison tests whether the extra structure explains generalizable variance rather than simply fitting more parameters.

\subsection{D3: Memory and plasticity endpoints}
Use a structured learning task with item identity, temporal order, source context, and relational composition. Perform encoding before the assigned sleep condition, retrieval after sleep, and delayed retrieval after recovery. Novel probes combine features not presented together during encoding. Attention, psychomotor vigilance, motivation, mood, and subjective sleepiness are measured so that alertness is not mistaken for consolidation.

For target memory $m$, define
\begin{equation}
\begin{aligned}
M(m)=1-\frac{1}{Z_m}\big(&E_{\mathrm{item}}(m)+E_{\mathrm{order}}(m)\\
&+E_{\mathrm{relation}}(m)+E_{\mathrm{source}}(m)\big),
\end{aligned}
\label{eq:memory}
\end{equation}
where $Z_m$ is a preregistered normalization and each error component is reported separately. Primary endpoints include item, order, relation, and source accuracy, confidence calibration, response time, interference resistance, and neural or physiological markers of reactivation where available. A stage feature is not promoted as causal unless it predicts held-out memory beyond alertness, encoding opportunity, practice, and nuisance controls and is supported by an intervention contrast.

\subsection{D4: Immune, maintenance, and systemic endpoints}
Use a panel rather than a single universal biomarker. Candidate immune measures include blood-cell subsets, inflammatory cytokines, symptom and infection surveillance, and an ethically approved standardized challenge where applicable. Candidate systemic measures include resting heart rate, heart-rate variability, blood pressure, glucose regulation proxies, mood, and endocrine markers. Clearance-related physiology is included only with validated, approved methods and is interpreted as a candidate mediator.

For participant $i$, endpoint $j$, and time $t$, use a hierarchical model:
\begin{equation}
Y_{i,j,t}=\theta_{0j}+\theta_{Hj}H_{i,t}+\theta_{Cj}C_{i,t}+\theta_{Z j}^{\top}z_{i,t}+b_{i,j}+\delta_{c,j}+\varepsilon_{i,j,t},
\label{eq:outcome}
\end{equation}
where $b_{i,j}$ is participant-specific baseline, $\delta_{c,j}$ is a condition effect, and $\varepsilon_{i,j,t}$ is residual variation. Domain-specific intervals and predictive checks are reported. A multi-marker immune movement is not translated into general “immune health” without a declared interpretation.

\subsection{D5: Causal mediation and utility}
The causal contrast for changing sleep feature $k$ on endpoint $j$ is
\begin{equation}
\tau_{k,j}=\mathbb{E}[Y_j\mid do(A_k=1),X]-\mathbb{E}[Y_j\mid do(A_k=0),X],
\label{eq:causal}
\end{equation}
where $A_k$ is a randomized or protocol-defined manipulation and $X$ includes preregistered covariates. Mediation by a stage or clearance measure is examined only if the mediator is reliable, temporally ordered, and measured without treatment-induced selection that defeats the interpretation.

An exploratory cross-domain utility layer is defined as
\begin{equation}
U_i=\sum_{j=1}^{J}w_{i,j}\widetilde{Y}_{i,j}-\lambda_i R_i,
\label{eq:utility}
\end{equation}
where $w_{i,j}$ are explicitly declared participant or clinical weights, $\widetilde{Y}_{i,j}$ are standardized outcome benefits, and $R_i$ is burden or safety risk. The weights cannot replace the scientific endpoint results. They make tradeoffs visible: a schedule that improves memory but worsens mood, immune markers, or accident risk is not a universally successful sleep intervention.

\subsection{D6: Individual response and decision gates}
Fit a hierarchical individual-response model with age, chronotype, baseline sleep, health, sex, medication, stress, and prior sleep as moderators. Cross-validation is performed by participant and session. The personalized controller selects only among safe, approved schedules and targets one declared endpoint at a time. It must outperform a duration-matched baseline on held-out data, preserve non-target outcomes within bounds, and remain calibrated for uncertainty.

A sleep feature is promoted as a candidate causal contributor only if it is reliably measured; separated from duration, timing, alertness, and prior wake; generalized to held-out sessions or participants; supported by a justified causal contrast; and accompanied by safety and non-target checks. If the duration-only model is equally predictive, the atlas simplifies its explanation. If memory improves while immune or safety endpoints deteriorate, the intervention fails the utility gate. Invalid assay batches, acute illness, missing stage labels, irregular medication, or sensor failure are marked `needs\_human\_input`, not converted into evidence.

\begin{table}[t]
\caption{SFCA promotion gates}
\label{tab:gates}
\centering
\footnotesize
\setlength{\tabcolsep}{2pt}
\begin{tabular}{p{0.24\columnwidth}p{0.61\columnwidth}}
\toprule
Gate & Required evidence \\
\midrule
Phenotype & Sleep feature and endpoint are defined before testing. \\
Separation & Effect survives duration, alertness, encoding, practice, and state controls. \\
Generalization & Result transfers to held-out sessions or participants. \\
Causality & Randomized or justified intervention contrast supports the direction. \\
Safety & Non-target cognitive, immune, mood, and safety outcomes remain bounded. \\
Personalization & Individual response is calibrated and improves a declared endpoint. \\
\bottomrule
\end{tabular}
\end{table}

\section{Expected Outcome Branches and Conditional Conclusions}
SFCA is designed to produce informative branches rather than one universal sleep score. If decomposed features outperform duration-only models, the result supports a multi-component sleep function. If duration alone performs equally well, the complexity of the proposed state representation is not justified for that endpoint. If prior wake and circadian phase separate under matched time in bed, the two-process model gains causal support. If they do not, the result narrows rather than destroys the usefulness of the operational distinction.

For memory, a selective improvement in relational or temporal-order retention with stable alertness would support a sleep-dependent plasticity or replay contribution. Improvement that disappears after controlling for encoding opportunity or vigilance would indicate a performance pathway instead. Preserved item accuracy with impaired source or relation accuracy would show that sleep protects some memory components more than others. A memory null result would not imply that sleep has no function; it would constrain the task and mechanism tested.

For immune and maintenance outcomes, a marker change with no cognitive consequence would establish a domain dissociation. A clearance-related measure that mediates an outcome under correct temporal ordering would support a candidate maintenance pathway, but not a sole purpose of sleep. A lack of mediation would leave other cellular, endocrine, and network mechanisms open. Because immune responses vary across people, a population average may hide a subgroup that benefits or is harmed by the same schedule.

A strong positive result would be an endpoint-specific, held-out, causal map: for example, a protected stage pattern predicts relational memory, a separate timing contrast predicts immune signaling, and a recovery schedule restores both without safety cost in a calibrated subgroup. Such a result would justify targeted sleep interventions and a richer theory of sleep need. It would not mean that every person needs the same stage distribution or that one night determines long-term health.

The near-term answer to the motivating question is direct. We need sleep because it coordinates several processes that waking cannot fully substitute: neural plasticity and memory processing, regulation of arousal and future learning capacity, cellular and fluid-transport maintenance, immune communication, and systemic repair. These functions are coupled but not identical. The practical scientific task is therefore to identify which component is needed for which outcome, at what time, and for whom.

\section{Risks, Limitations, and Human Review Requirements}
\subsection{Scientific and technical risks}
The largest risk is confounding. Sleep loss changes alertness, mood, motivation, encoding opportunity, hormonal state, and task strategy simultaneously. Circadian displacement changes light, meal, and social schedules. Repeated memory testing produces practice. Blood sampling and imaging introduce burden and may change sleep. These risks are addressed through counterbalancing, washout, fixed schedules, vigilance controls, novelty probes, repeated-measures models, and preregistered missing-data rules.

Measurement validity is another risk. Consumer wearables may estimate duration but not reliably distinguish stages. Polysomnography provides richer staging but is intrusive. Inflammatory markers fluctuate and can be influenced by minor illness. Clearance-related measurements may be indirect and difficult to compare across species. SFCA therefore reports reliability, modality, uncertainty, and domain-specific endpoints rather than treating one sensor or biomarker as ground truth.

The model can also overfit individual differences. A hierarchical controller may learn who performs well under a particular schedule without identifying why. Participant and session holdouts, calibration, prespecified moderators, and a duration-matched baseline limit this risk. A causal-looking latent state is not accepted when its predictions are explained by leakage, subject identity, arousal, or measurement quality.

\subsection{Safety, ethics, and governance}
Sleep restriction and circadian displacement can worsen mood, attention, seizure susceptibility, metabolic regulation, and accident risk in vulnerable individuals. Clinical and psychiatric screening, stop rules, recovery access, monitoring, compensation, and post-study support are mandatory. Participants must understand the schedule, risks, data collection, and right to withdraw without penalty. No unsupervised sleep-loss prescription follows from this paper.

Longitudinal wearable, neural, endocrine, and immune data can reveal sensitive health patterns. Data minimization, de-identification, access control, retention limits, and participant governance are required. Immune challenges, blood collection, imaging, and any future neuromodulation need separate institutional review. A proposed benefit to memory does not override privacy, autonomy, safety, or equitable recruitment.

\subsection{Human review and reproducibility}
Before execution, experts must inspect sample size and power, carryover assumptions, sleep-disorder exclusions, stage-scoring reliability, immune assay validity, clearance measurement, medication and illness handling, causal mediation assumptions, and personalized-controller calibration. The final dataset must retain condition order, washout, raw and processed sleep features, model versions, random seeds, fold assignments, protocol deviations, adverse events, null results, and all unknowns. Failed batches remain failure records and are marked `needs\_human\_input`.

\section{Conclusion}
The question “Why do we need sleep?” has a richer answer than a single-purpose slogan. Sleep is a regulated state in which homeostatic pressure, circadian timing, continuity, stage dynamics, and individual physiology coordinate multiple forms of neural and bodily maintenance. Memory processing, plasticity, immune communication, clearance-related physiology, and systemic repair can share pathways while retaining distinct causal signatures.

The Sleep Function Causal Atlas makes this answer experimentally actionable. It separates the two-process architecture, measures sleep as a vector, evaluates content-sensitive memory and multi-marker physiological outcomes, and uses safe factorial contrasts, recovery, held-out generalization, and individual-response gates. Its strongest positive result would be a selective map of which sleep feature protects which endpoint and for whom. Its strongest negative result would show that a proposed feature adds no predictive or causal value beyond duration and alertness. Both outcomes improve the science.

The present artifact contains no observed intervention. It establishes a path from evidence to a research design that is ambitious enough to discover multiple functions of sleep and disciplined enough to distinguish association from causation. We need sleep not because one explanation has already won, but because the living system repeatedly uses this regulated state to maintain capabilities that waking alone does not reliably preserve.

\appendices
\section{Sleep Phenotype and Decision Rules Appendix}
Table~\ref{tab:phenotype} records the phenotype fields that the Author Agent must preserve. No field is a substitute for another, and no composite utility score can erase a domain-specific adverse effect.

\begin{table}[h]
\caption{Operational sleep phenotype}
\label{tab:phenotype}
\centering
\footnotesize
\setlength{\tabcolsep}{2pt}
\begin{tabular}{p{0.24\columnwidth}p{0.59\columnwidth}}
\toprule
Field & Operational role \\
\midrule
Duration & Time asleep and time in bed, reported separately. \\
Timing & Clock phase, light exposure, and circadian proxy. \\
Continuity & Efficiency, awakenings, and fragmentation. \\
Stage structure & Slow-wave activity, spindles, REM timing, and transitions. \\
Prior wake & Accumulated wake and recovery history. \\
Memory & Item, order, source, relational, confidence, and interference endpoints. \\
Physiology & Immune, maintenance, cardiovascular, metabolic, endocrine, and mood panels. \\
Individual state & Age, chronotype, health, medication, stress, and baseline response. \\
Safety & Sleepiness, adverse events, mood, cognition, and operational risk. \\
\bottomrule
\end{tabular}
\end{table}

The promotion decision is conjunctive: a finding must have a defined phenotype, reliable measurement, confound separation, held-out generalization, a justified causal contrast, and acceptable non-target and safety outcomes. A prediction can remain useful at a lower level without being promoted to a mechanistic or clinical claim.

\section{Evidence Coverage and Review Appendix}
The Survey-to-Author binding uses project identifier \texttt{neur\_4}, Survey run identifier \texttt{neur\_4-survey-001}, and project-context fingerprint \texttt{neur4-sleep-function-context-v1}. Canonical upstream artifacts are the Survey evidence plan, gap ledger, Idea result, ExperimentDesign JSON, and Author handoff stored with this report. Accepted gaps are G1--G7; G8 and G9 remain candidate gaps requiring validation. The source registry is E1--E9.

Claims about sleep regulation and health effects are bounded by E1 and E2. Claims about memory and plasticity are bounded by E3 and E4. Clearance-related claims are bounded by E5 and E6. Immune claims are bounded by E7 and E8. Individual variability claims are bounded by E9. The document status is \texttt{PROPOSAL\_NO\_OBSERVED\_RESULTS}; no sleep manipulation, assay, or clinical outcome was executed.

\begin{thebibliography}{9}
\bibitem{NHLBI2022} National Heart, Lung, and Blood Institute, ``Sleep deprivation and deficiency,'' U.S. National Institutes of Health, updated Mar. 24, 2022. [Online]. Available: \url{https://www.nhlbi.nih.gov/health/sleep-deprivation}
\bibitem{BorbeLy1982} A. A. Borb\'ely, ``A two process model of sleep regulation,'' \textit{Human Neurobiology}, vol. 1, no. 3, pp. 195--204, 1982.
\bibitem{Walker2006} M. P. Walker and R. Stickgold, ``Sleep, memory, and plasticity,'' \textit{Annual Review of Psychology}, vol. 57, pp. 139--166, 2006.
\bibitem{RaschBorn2013} B. Rasch and J. Born, ``About sleep's role in memory,'' \textit{Physiological Reviews}, vol. 93, no. 2, pp. 681--766, 2013.
\bibitem{Klinzing2019} J. G. Klinzing, N. Niethard, and J. Born, ``Mechanisms of systems memory consolidation during sleep,'' \textit{Nature Neuroscience}, vol. 22, pp. 1598--1610, 2019.
\bibitem{Xie2013} L. Xie \textit{et al.}, ``Sleep drives metabolite clearance from the adult brain,'' \textit{Science}, vol. 342, no. 6156, pp. 373--377, 2013.
\bibitem{Nedergaard2020} M. Nedergaard and S. A. Goldman, ``Glymphatic failure as a final common pathway to dementia,'' \textit{Science}, vol. 370, no. 6521, pp. 50--56, 2020.
\bibitem{Besedovsky2019} L. Besedovsky, T. Lange, and M. Haack, ``The sleep-immune crosstalk in health and disease,'' \textit{Physiological Reviews}, vol. 99, no. 3, pp. 1325--1380, 2019.
\bibitem{Irwin2015} M. R. Irwin, ``Why sleep is important for health: A psychoneuroimmunology perspective,'' \textit{Annual Review of Psychology}, vol. 66, pp. 143--172, 2015.
\end{thebibliography}

\end{document}
